% Aufgabensammlung — Kapitel 2
% Das Kapitel wird durch den Namen dieser Datei bestimmt.

\begin{exo}[Von kaltem Eis zu siedendem Wasser: Größenordnungen]{chauffage-glace-eau-ebullition}
\seoready{true}
\keywords{Kalorimetrie, Phasenübergang, Schmelzen, Verdampfen, latente Wärme, Größenordnung}

\textbf{Zahlenwerte (bei Atmosphärendruck):}
\[
\begin{aligned}
c_{\text{Eis}} &= 2{,}1\times10^3\,\text{J}\!\cdot\!\text{kg}^{-1}\!\cdot\!\text{K}^{-1},
& L_f &= 3{,}34\times10^5\,\text{J}\!\cdot\!\text{kg}^{-1},\\
c_{\text{Wasser}} &= 4{,}18\times10^3\,\text{J}\!\cdot\!\text{kg}^{-1}\!\cdot\!\text{K}^{-1},
& L_v &= 2{,}26\times10^6\,\text{J}\!\cdot\!\text{kg}^{-1}.
\end{aligned}
\]
Bei Atmosphärendruck wird ein Kilogramm Eis, das anfangs eine Temperatur von $-100\,{}^\circ\text{C}$ hat, schrittweise erwärmt.

\begin{enumerate}
  \item Berechnen Sie die Energie, die erforderlich ist, um das Eis von $-100\,{}^\circ\text{C}$ auf $0\,{}^\circ\text{C}$ zu erwärmen.
  \item Berechnen Sie die Energie, die erforderlich ist, um das Eis bei $0\,{}^\circ\text{C}$ vollständig zu schmelzen.
  \item Berechnen Sie die Energie, die erforderlich ist, um das flüssige Wasser von $0\,{}^\circ\text{C}$ auf $100\,{}^\circ\text{C}$ zu erwärmen.
  \item Berechnen Sie die zusätzliche Energie, die erforderlich ist, um das Wasser bei $100\,{}^\circ\text{C}$ vollständig zu verdampfen.
  \item Welcher Schritt erfordert die meiste Energie? Vergleichen Sie die Energie zum Erwärmen eines Kilogramms Wasser von $0\,{}^\circ\text{C}$ auf $100\,{}^\circ\text{C}$ mit der potenziellen Energie einer Masse $m$, die zehn Meter tief fällt. Welche Masse $m$ müsste fallen, um dieselbe Energie freizusetzen?
  \item Betrachten Sie nun den umgekehrten Vorgang. Eine Masse $m_v$ Wasserdampf kondensiert auf einer Windschutzscheibe, ohne dass das entstandene flüssige Wasser anschließend abkühlt. Geben Sie die bei der Kondensation freigesetzte Wärme $Q_{\text{frei}}$ an und berechnen Sie sie für $m_v=1{,}0\times10^{-2}\,\text{kg}$. Bei der Temperatur der Windschutzscheibe sei $L_v=2{,}45\times10^6\,\text{J}\!\cdot\!\text{kg}^{-1}$.
\end{enumerate}

\begin{indication}
Zum Erwärmen einer Phase ohne Phasenübergang gilt $Q=mc\Delta T$. Für einen Phasenübergang bei konstanter Temperatur gilt $Q=mL$. Die potenzielle Energie einer Masse $m$ in der Höhe $h$ ist $E_p=mgh$; verwenden Sie $g=9{,}81\,\text{m}\!\cdot\!\text{s}^{-2}$.
\end{indication}

\begin{solution}
\textbf{1.} Das Erwärmen des Eises bis zu seiner Schmelztemperatur erfordert
\[
Q_1=mc_{\text{Eis}}\Delta T
=1{,}0\times2{,}1\times10^3\times100
=2{,}1\times10^5\,\text{J}.
\]

\textbf{2.} Während des Schmelzens bleibt die Temperatur bei $0\,{}^\circ\text{C}$:
\[
Q_2=mL_f=1{,}0\times3{,}34\times10^5
=3{,}34\times10^5\,\text{J}.
\]

\textbf{3.} Zum Erwärmen des flüssigen Wassers von $0$ auf $100\,{}^\circ\text{C}$ gilt
\[
Q_3=mc_{\text{Wasser}}\Delta T
=1{,}0\times4{,}18\times10^3\times100
=4{,}18\times10^5\,\text{J}.
\]

\textbf{4.} Die vollständige Verdampfung erfordert zusätzlich
\[
Q_4=mL_v=1{,}0\times2{,}26\times10^6
=2{,}26\times10^6\,\text{J}.
\]
Allein dieser Schritt benötigt also eine Energie von einigen Megajoule.

\textbf{5.} Die Verdampfung ist mit Abstand der energieaufwendigste Schritt:
\[
Q_4=2{,}26\times10^6\,\text{J}
>Q_3=4{,}18\times10^5\,\text{J}
>Q_2=3{,}34\times10^5\,\text{J}
>Q_1=2{,}1\times10^5\,\text{J}.
\]
Insbesondere erfordert die Verdampfung etwa $(2{,}26\times10^6)/(4{,}18\times10^5)\approx5{,}4$-mal so viel Energie wie das Erwärmen des flüssigen Wassers von $0$ auf $100\,{}^\circ\text{C}$.

Für eine Masse $m$, die aus einer Höhe $h=10\,\text{m}$ fällt, gilt $E_p=mgh$. Aus $mgh=Q_3$ folgt
\[
m=\frac{Q_3}{gh}
=\frac{4{,}18\times10^5}{9{,}81\times10}
\approx4{,}26\times10^3\,\text{kg}.
\]
Eine Masse von ungefähr $4{,}3$ Tonnen müsste also zehn Meter tief fallen, um ebenso viel Energie freizusetzen, wie zum Erwärmen eines Kilogramms Wasser von $0$ auf $100\,{}^\circ\text{C}$ nötig ist! Das ist enorm und erklärt, warum die Warmwasserbereitung einen bedeutenden Anteil des häuslichen Energieverbrauchs ausmacht. Zudem besitzt Wasser eine sehr große spezifische Wärmekapazität, insbesondere im Vergleich zu üblichen Metallen; sie ist beispielsweise etwa zehnmal so groß wie die von Kupfer (siehe die folgenden Aufgaben).

\textbf{6.} Die Kondensation ist der umgekehrte Vorgang der Verdampfung: Der Dampf gibt die latente Wärme an die Windschutzscheibe ab,
\[
Q_{\text{frei}}=m_vL_v.
\]
Für $m_v=1{,}0\times10^{-2}\,\text{kg}$ ergibt sich
\[
Q_{\text{frei}}=1{,}0\times10^{-2}\times2{,}45\times10^6
=2{,}45\times10^4\,\text{J}.
\]
Somit kann bereits die Kondensation einer kleinen Dampfmasse eine nicht zu vernachlässigende Wärmemenge freisetzen. Diese Wärme wird von der Windschutzscheibe und ihrer Umgebung aufgenommen.
\end{solution}
\end{exo}
\begin{exo}[Evapotranspiration eines großen Baumes: eine natürliche Klimaanlage?]{odg-evapotranspiration-arbre}
\seoready{true}
\keywords{Größenordnung, Evapotranspiration, Baum, latente Wärme, Kühlleistung, Klimaanlage, COP}

An einem heißen und trockenen Tag kann ein großer Baum bei ausreichender Wasserversorgung ungefähr $500\,\text{L}=5{,}0\times10^{-1}\,\text{m}^3$ Wasser evapotranspirieren. Da die Transpiration hauptsächlich bei Licht stattfindet, nehmen wir an, dass sie sich über zwölf Stunden verteilt. Die Baumkrone bedecke am Boden eine Fläche $A_{\text{Baum}}=50\,\text{m}^2$. Gegeben sind die Dichte des Wassers und seine Verdampfungsenthalpie bei Atmosphärendruck:
\[
\rho_{\text{Wasser}}=1{,}0\times10^3\,\text{kg}\!\cdot\!\text{m}^{-3},
\qquad L_v=2{,}45\times10^6\,\text{J}\!\cdot\!\text{kg}^{-1}.
\]

\begin{enumerate}
  \item Berechnen Sie die während des Tages evapotranspirierte Wassermasse und die zu ihrer Verdampfung erforderliche Energie.
  \item Erklären Sie, warum dadurch eine Kühlwirkung entsteht.
  \item Berechnen Sie während der zwölf Stunden aktiver Transpiration die mittlere Kühlleistung des Baumes pro Quadratmeter.
  \item Zum Vergleich betrachten wir eine Klimaanlage mit elektrischer Leistungsaufnahme $P_{\mathrm{el}}=2{,}0\times10^3\,\text{W}$, einer Leistungszahl $\mathrm{COP}=3{,}0$ und einer gekühlten Raumfläche $A_{\mathrm{Raum}}=20\,\text{m}^2$. Definitionsgemäß gilt $\mathrm{COP}=P_{\mathrm{cool}}/P_{\mathrm{el}}$. Berechnen Sie die Kühlleistung und die Kühlleistung pro Quadratmeter und vergleichen Sie diese Werte mit denen des Baumes.
  \item Warum kann man nicht folgern, dass Baum und Klimaanlage genau dieselbe empfundene Kühlung bewirken, selbst wenn ihre Leistungen pro Flächeneinheit dieselbe Größenordnung haben?
  \item Kann man eine Wohnung mit vielen Zimmerpflanzen kühlen? (Allgemeinwissen: Die Antwort hängt von biologischen Vorgängen ab.)
\end{enumerate}

\begin{indication}
Verwenden Sie $m=\rho V$, $Q_{\mathrm{verd}}=mL_v$ und $P=Q/\tau$. Für die Klimaanlage gilt $P_{\mathrm{cool}}=\mathrm{COP}\,P_{\mathrm{el}}$.
\end{indication}

\begin{solution}
\textbf{1.} Die entsprechende Wassermasse ist
\[
m=\rho_{\text{Wasser}}V
=1{,}0\times10^3\times5{,}0\times10^{-1}
=5{,}0\times10^2\,\text{kg}.
\]
Die zur Verdampfung erforderliche Energie beträgt daher
\[
Q_{\mathrm{verd}}=mL_v
=5{,}0\times10^2\times2{,}45\times10^6
=1{,}225\times10^9\,\text{J}.
\]
Die Größenordnung ist also ein Gigajoule.

\textbf{2.} Die Verdampfung ist ein endothermer Phasenübergang: Die Energie wird den Blättern, dem Boden und der umgebenden Luft entzogen. Dieser Wärmeentzug senkt die Temperatur der Blätter und ihrer unmittelbaren Umgebung und erzeugt so die Kühlwirkung.

\textbf{3.} Zwölf Stunden entsprechen $\tau=12\times3600=4{,}32\times10^4\,\text{s}$. Damit beträgt die mittlere Kühlleistung
\[
P_{\text{Baum}}=\frac{Q_{\mathrm{verd}}}{\tau}
=\frac{1{,}225\times10^9}{4{,}32\times10^4}
\approx2{,}84\times10^4\,\text{W},
\]
und bezogen auf die überdeckte Fläche
\[
\frac{P_{\text{Baum}}}{A_{\text{Baum}}}
=\frac{2{,}84\times10^4}{50}
\approx5{,}7\times10^2\,\text{W}\!\cdot\!\text{m}^{-2}.
\]

\textbf{4.} Die nutzbare Kühlleistung der Klimaanlage ist
\[
P_{\mathrm{cool}}=\mathrm{COP}\,P_{\mathrm{el}}
=3{,}0\times2{,}0\times10^3
=6{,}0\times10^3\,\text{W}.
\]
Bezogen auf die Raumfläche gilt
\[
\frac{P_{\mathrm{cool}}}{A_{\mathrm{Raum}}}
=\frac{6{,}0\times10^3}{20}
=3{,}0\times10^2\,\text{W}\!\cdot\!\text{m}^{-2}.
\]
In diesem Modell ist die Evapotranspirationsleistung des Baumes pro Flächeneinheit somit fast doppelt so groß wie die einer üblichen Klimaanlage.

\textbf{5.} Der Vergleich betrifft nur Energieflüsse. Die Klimaanlage kühlt einen geschlossenen und kontrollierten Raum, während sich der vom Baum erzeugte Wasserdampf in der Atmosphäre verteilt und der Wind warme Luft heranführt. Die empfundene Wirkung hängt daher von Lüftung, Luftfeuchte, Umgebungstemperatur und Abstand zum Baum ab. Außerdem variiert die Transpiration mit Sonneneinstrahlung, Wind, Luftfeuchte, Baumart und Wasserangebot und kann bei Trockenheit stark abnehmen. Andererseits spendet der Baum Schatten und verringert dadurch unmittelbar die auf Boden und Menschen treffende Sonnenstrahlung. Die berechneten Leistungen ermöglichen also einen Vergleich der Größenordnungen, aber keine exakte Gleichsetzung des thermischen Komforts.

\textbf{6.} In der Praxis wäre die Kühlwirkung von Zimmerpflanzen sehr gering. Bei den meisten Pflanzen fördert Licht die Öffnung der Stomata, kleiner Poren in den Blättern, durch die Wasserdampf entweicht. Innenräume sind meist schwach beleuchtet, sodass sich die Stomata weniger öffnen und die Transpiration abnimmt. In einem schlecht belüfteten Raum bremst zudem die zunehmende Luftfeuchte allmählich die Verdampfung.

Viele Pflanzen können daher eine leichte lokale Verdunstungskühlung bewirken, ersetzen aber keine Klimaanlage. Um dem Gebäude dauerhaft Energie zu entziehen, müsste die feuchte Luft nach außen abgeführt werden. Eine intensive künstliche Beleuchtung könnte die Transpiration anregen, doch ihre elektrische Energie würde schließlich fast vollständig als Wärme im Raum enden. An trockene Lebensräume angepasste CAM-Pflanzen, darunter Kakteen, bilden eine Ausnahme: Ihre Stomata öffnen sich hauptsächlich nachts.
\end{solution}
\end{exo}

\begin{exo}[Mischen zweier Wassermassen]{calorimetrie-melange-eau}
\seoready{true}
\keywords{Kalorimetrie, Mischung, Wärmekapazität, spezifische Wärmekapazität, Joseph Black, Gleichgewichtstemperatur}

$m_1=0{,}200\,\text{kg}$ Wasser bei $T_1=80\,^\circ\text{C}$ werden in ein Gefäß gegossen, das bereits $m_2=0{,}300\,\text{kg}$ Wasser bei $T_2=15\,^\circ\text{C}$ enthält. Das Gefäß ist ein ideales Kalorimeter: Wärmeverluste an die Umgebung sowie seine eigene Wärmekapazität werden vernachlässigt. Es gilt
\[
Q=mc\Delta T,
\]
wobei $c=4{,}18\times10^3\,\text{J}\!\cdot\!\text{kg}^{-1}\!\cdot\!\text{K}^{-1}$ die spezifische Wärmekapazität von flüssigem Wasser ist.

\begin{enumerate}
  \item Begründen Sie, dass die vom warmen Wasser abgegebene Wärme vollständig vom kalten Wasser aufgenommen wird. Formulieren Sie die Erhaltungsgleichung mit $m_1$, $m_2$, $c$, $T_1$, $T_2$ und der Gleichgewichtstemperatur $T_{eq}$.
  \item Leiten Sie einen Ausdruck für $T_{eq}$ her und berechnen Sie den Zahlenwert.
  \item Berechnen Sie die beim Mischen vom warmen Wasser abgegebene Wärmemenge $Q$.
  \item Lässt sich mit solchen Mischungen derselben Substanz der Wert von $c$ bestimmen? Begründen Sie.
\end{enumerate}

\begin{indication}
Da das Kalorimeter isoliert und starr ist, bleibt die gesamte innere Energie $U_1+U_2$ erhalten: Die von einem Teil abgegebene Wärme wird vom anderen mit umgekehrtem Vorzeichen aufgenommen.
\end{indication}

\begin{solution}
\textbf{1.} Die Erhaltung der inneren Energie des Systems ergibt
\[
m_1c(T_{eq}-T_1)+m_2c(T_{eq}-T_2)=0.
\]

\textbf{2.} Der gemeinsame Faktor $c$ kürzt sich heraus:
\[
T_{eq}=\frac{m_1T_1+m_2T_2}{m_1+m_2}
=\frac{0{,}200\times80+0{,}300\times15}{0{,}200+0{,}300}
=41\,^\circ\text{C}.
\]

\textbf{3.} Das warme Wasser gibt ab
\[
Q=m_1c(T_1-T_{eq})
=0{,}200\times4{,}18\times10^3\times(80-41)
\approx3{,}26\times10^4\,\text{J}.
\]
Das kalte Wasser nimmt genau dieselbe Wärmemenge auf:
\[
m_2c(T_{eq}-T_2)
=0{,}300\times4{,}18\times10^3\times26
\approx3{,}26\times10^4\,\text{J}.
\]

\textbf{4.} Nein. Bei zwei Massen derselben Substanz multipliziert derselbe Faktor $c$ beide Terme der Bilanz und kürzt sich heraus. Die gemessene Gleichgewichtstemperatur hängt daher nicht von $c$ ab, sondern ist nur das mit den Massen gewichtete Mittel der Anfangstemperaturen. Zur Bestimmung einer spezifischen Wärmekapazität durch die Mischungsmethode benötigt man eine Substanz mit bekannter Wärmekapazität, wie in der folgenden Aufgabe.
\end{solution}
\end{exo}

\begin{exo}[Bestimmung der spezifischen Wärmekapazität eines Metalls mit der Mischungsmethode]{calorimetrie-methode-melanges}
\seoready{true}
\keywords{Kalorimetrie, Mischungsmethode, spezifische Wärmekapazität, Kalorimeter, Wasserwert}

Wie bereits Joseph Black im 18. Jahrhundert soll die spezifische Wärmekapazität $c_m$ eines unbekannten Metalls mit der \emph{Mischungsmethode} bestimmt werden. Eine Probe der Masse $m=0{,}150\,\text{kg}$ wird auf $T_m=95\,^\circ\text{C}$ erhitzt und rasch in ein Kalorimeter mit $m_e=0{,}250\,\text{kg}$ Wasser getaucht. Für das Wasser gilt $c_e=4{,}18\times10^3\,\text{J}\!\cdot\!\text{kg}^{-1}\!\cdot\!\text{K}^{-1}$, seine Anfangstemperatur ist $T_e=18\,^\circ\text{C}$. Nach Erreichen des thermischen Gleichgewichts wird $T_{eq}=22\,^\circ\text{C}$ gemessen. Zunächst wird die Wärmekapazität von Gefäß, Rührer und Thermometer vernachlässigt.

\begin{enumerate}
  \item Stellen Sie die Energieerhaltung für das isolierte System aus Metall und Wasser auf, wobei $c_m$ die einzige Unbekannte ist.
  \item Leiten Sie einen Ausdruck für $c_m$ her und berechnen Sie den Zahlenwert. Welchem gebräuchlichen Metall entspricht er ungefähr? Vergleichswerte in $\text{J}\!\cdot\!\text{kg}^{-1}\!\cdot\!\text{K}^{-1}$: Aluminium $9{,}0\times10^2$, Eisen $4{,}5\times10^2$, Kupfer $3{,}9\times10^2$, Blei $1{,}3\times10^2$.
  \item Das Kalorimetergefäß nimmt tatsächlich einen Teil der Wärme auf. Dieser Effekt werde durch den \emph{Wasserwert} $\mu=1{,}5\times10^{-2}\,\text{kg}$ modelliert: Das Kalorimeter verhält sich wie eine zusätzliche Wassermasse $\mu$. Berechnen Sie $c_m$ erneut und erläutern Sie die Richtung der Korrektur.
  \item Warum muss die Probe rasch eingetaucht und das Kalorimeter während der Messung gut gegen die Umgebungsluft isoliert werden?
\end{enumerate}

\begin{indication}
Das Metall gibt $Q_m=mc_m(T_m-T_{eq})$ ab. Diese Wärme wird vollständig vom Wasser und in Aufgabe 3 zusätzlich vom Kalorimeter aufgenommen: $Q_m=(m_e+\mu)c_e(T_{eq}-T_e)$.
\end{indication}

\begin{solution}
\textbf{1.} Für das isolierte System gilt
\[
mc_m(T_m-T_{eq})=m_ec_e(T_{eq}-T_e).
\]

\textbf{2.} Auflösen nach $c_m$ liefert
\[
c_m=\frac{m_ec_e(T_{eq}-T_e)}{m(T_m-T_{eq})}
=\frac{0{,}250\times4{,}18\times10^3\times(22-18)}{0{,}150\times(95-22)}
\approx3{,}82\times10^2\,\text{J}\!\cdot\!\text{kg}^{-1}\!\cdot\!\text{K}^{-1}.
\]
Dieser Wert liegt sehr nahe bei dem von Kupfer.

\textbf{3.} Mit dem Wasserwert des Kalorimeters ergibt sich
\[
c_m=\frac{(m_e+\mu)c_e(T_{eq}-T_e)}{m(T_m-T_{eq})}
=\frac{(0{,}250+0{,}015)\times4{,}18\times10^3\times4}{0{,}150\times73}
\approx4{,}05\times10^2\,\text{J}\!\cdot\!\text{kg}^{-1}\!\cdot\!\text{K}^{-1}.
\]
Die Korrektur erhöht $c_m$: Ohne das Gefäß wurde die tatsächlich vom Metall abgegebene Wärme und damit seine spezifische Wärmekapazität unterschätzt.

\textbf{4.} Die Methode setzt während der gesamten Messung ein isoliertes System voraus. Rasches Eintauchen begrenzt den Wärmeaustausch der heißen Probe mit der Luft, bevor sie das Wasser erreicht; eine gute Isolierung begrenzt Verluste, während sich das Gleichgewicht einstellt. Jeder nicht berücksichtigte Wärmestrom verfälscht die Bilanz und den daraus bestimmten Wert von $c_m$. Genau diese experimentelle Sorgfalt zeigten bereits Black sowie später Lavoisier und Laplace mit ihrem Eiskalorimeter.
\end{solution}
\end{exo}

\begin{exo}[Nahrungsmittelkalorien und Stoffwechselleistung]{calorimetrie-calories-alimentaires}
\seoready{true}
\keywords{Kalorie, Kilokalorie, Nahrungsmittelkalorie, mechanisches Wärmeäquivalent, Stoffwechselleistung, Einheiten}

Eine Nährwertangabe geht davon aus, dass ein Erwachsener im Mittel etwa $2000\,\text{kcal}$ pro Tag aufnimmt. Die Kilokalorie ist eine in der Ernährung noch verwendete Nicht-SI-Einheit; ihre Umrechnung in SI-Einheiten lautet $1\,\text{kcal}=4{,}18\times10^3\,\text{J}$.

\begin{enumerate}
  \item Rechnen Sie diese tägliche Energiezufuhr in Joule um.
  \item Bestimmen Sie daraus die mittlere Stoffwechselleistung in Watt, wenn die Energie gleichmäßig über 24 Stunden umgesetzt wird.
  \item Ein Energieriegel trägt die Angabe „$210\,\text{kcal}$“. Welche Masse Wasser könnte mit derselben Energiemenge von anfangs $20\,^\circ\text{C}$ bis zum Siedepunkt bei $100\,^\circ\text{C}$ erwärmt werden, wenn sie vollständig in Wärme umgewandelt würde? Es gilt $c_{\text{Wasser}}=4{,}18\times10^3\,\text{J}\!\cdot\!\text{kg}^{-1}\!\cdot\!\text{K}^{-1}$.
  \item In welcher Form wird diese Energie vom menschlichen Körper tatsächlich umgesetzt? Antworten Sie qualitativ.
\end{enumerate}

\begin{indication}
Verwenden Sie für Aufgabe 2 $\mathcal P=E/\Delta t$. Rechnen Sie für Aufgabe 3 zunächst die Energie in Joule um und lösen Sie dann $Q=mc\Delta T$ nach $m$ auf.
\end{indication}

\begin{solution}
\textbf{1.} In SI-Einheiten gilt
\[
E=2000\times4{,}18\times10^3
=8{,}36\times10^6\,\text{J}.
\]

\textbf{2.} Die mittlere Leistung beträgt
\[
\mathcal P=\frac{8{,}36\times10^6}{8{,}64\times10^4}\approx97\,\text{W}.
\]
Über den ganzen Tag setzt ein Mensch also im Mittel ungefähr dieselbe Leistung um wie eine Glühlampe. Diese oft überraschende Größenordnung zeigt, wie nützlich die Umrechnung von Nahrungsmittelkalorien in vertraute physikalische Einheiten ist.

\textbf{3.} Die Energie des Riegels ist
\[
Q=210\times4{,}18\times10^3\approx8{,}78\times10^5\,\text{J}.
\]
Mit $\Delta T=80\,\text{K}$ folgt
\[
m=\frac{Q}{c\Delta T}
=\frac{8{,}78\times10^5}{4{,}18\times10^3\times80}
\approx2{,}63\,\text{kg}.
\]
Ein einziger Energieriegel enthält größenordnungsmäßig also genug Energie, um etwa $2{,}5\,\text{kg}$ Leitungswasser zum Sieden zu bringen.

\textbf{4.} Der Körper „verbraucht“ Energie nicht in dem Sinn, dass sie verschwindet, sondern wandelt die chemische Energie der Nährstoffe um. Ein Teil dient der Muskelarbeit, der Funktion der Organe, der Aufrechterhaltung elektrischer Zellgradienten, dem Stofftransport und chemischen Synthesen. Ein weiterer Teil kann als Glykogen oder Fett chemisch gespeichert werden. Letztlich wird der größte Teil der umgesetzten Energie als Wärme dissipiert; diese trägt zur Aufrechterhaltung der Körpertemperatur bei und wird danach an die Umgebung abgegeben. Ein kleiner Anteil verlässt den Körper zudem als chemische Energie in Ausscheidungsprodukten.
\end{solution}
\end{exo}
